\chapter{Related Work}
\label{chapter_related_work}

UAV-assisted disaster response has been studied from several directions:
perception and scene understanding, survivor localization, cooperative search,
delivery logistics, resource-aware mission planning, and adaptive UAV
coordination. These works provide many of the upstream capabilities needed in a
post-earthquake response system. The focus of this thesis is narrower and
downstream: once survivor information has been accepted by a command center,
PGBM plans how responder UAVs should deliver temporary assistance before ground
SAR teams arrive.

\section{UAV-Based Disaster Perception and Search}
\label{sec_related_detection}

One major stream of research studies how UAVs can observe disaster scenes and
identify useful response information. Jankovic et al. use an optimized Swin
Transformer for UAV-assisted disaster-scene classification
\cite{jankovic2025transformer}. Their work is relevant because it shows how
aerial images can support emergency assessment, although the main task is scene
classification rather than survivor-assistance planning.

Several works move closer to survivor localization. ZaidAlkilani et al.
combine multi-UAV coverage with YOLOv8X-based human detection and geotagged
survivor reporting \cite{zaidalkilani2026aiuav}. Che et al. study
UAV-assisted 5G-NR localization using time-of-flight, angle-of-arrival, and
received-signal-strength measurements in complex disaster scenarios
\cite{che2025uav5g}. G\"uzey et al. formulate post-disaster victim
localization as a cooperative multi-agent problem with energy and
communication constraints \cite{guzey2026fedqmix}. Together, these studies
show that survivor information can be obtained through visual, radio-based, or
cooperative sensing pipelines.

Scene understanding is also important because a detected survivor location may
not be directly reachable by a UAV. Sirma et al. introduce D'RespNeT, a
post-earthquake UAV dataset and YOLOv8-based instance-segmentation model for
detecting collapsed structures, debris, civilians, access openings, and safe
UAV landing locations \cite{sirma2025drespnet}. Such information is useful for
choosing accessible drop-off points and feasible operating regions. More
generally, low-altitude disaster imagery datasets such as LADI and AIDER
support disaster-scene interpretation from aerial images
\cite{scheele2025ladi,kyrkou2020emergencynet}.

Another related stream focuses on how scout UAVs should search the affected
environment. Aminzadeh and Khoshnood propose a heterogeneous multi-UAV search
architecture for post-disaster assessment \cite{aminzadeh2023multiuav}, while
Kareem et al. use bio-inspired swarm coordination for thermal survivor search
\cite{kareem2026bioinspired}. Leong et al. consider dynamic decentralized 3D
urban coverage, where UAVs coordinate over three-dimensional observation
viewpoints \cite{leong2024urban3d}. These studies are important for the
upstream scout-UAV layer assumed in this thesis.

Overall, the perception, localization, and search literature mainly addresses
how survivor or environment information can be obtained. PGBM does not replace
these methods. Instead, it assumes that the command center has already accepted
survivor information and then asks a different question: how should responder
UAVs provide temporary assistance after detection but before physical rescue?

\section{UAV Delivery, Resource-Aware and Adaptive Mission Planning}
\label{sec_related_resource}

UAV delivery and mission-planning studies show that task decisions cannot be
separated from vehicle limitations. A delivery UAV has finite battery energy,
payload capacity, flight time, and route options. These constraints become more
important in disaster response because the environment may contain damaged
structures, blocked areas, and uncertain access conditions.

Delivery-oriented studies make this point directly. Dorling et al. include
payload weight, battery limits, delivery-time constraints, and drone reuse in
drone-specific routing models \cite{dorling2017dronevrp}. Cheng et al. model
payload-dependent energy consumption in multi-trip drone routing
\cite{cheng2020droneenergy}. These works are not written for post-earthquake
survivor assistance, but they show why delivery feasibility cannot be separated
from carried load and available energy.

Disaster logistics adds another layer to the routing problem. Tureci-Isik et
al. study humanitarian aid delivery using parallel truck-drone operations under
road-network disruption and uncertain travel times
\cite{tureciisik2026truckdrone}. Liu et al. and Xiong et al. consider UAV
mission assignment and path planning for disaster rescue, with Xiong et al.
placing particular emphasis on 3D path planning
\cite{liu2021multiuv,xiong2023multidrone}. These studies are useful for PGBM
because the responder UAVs must also plan feasible routes in a damaged
environment.

Resource-aware emergency-response planning has also been studied beyond
physical delivery. Akter et al. consider task offloading, trajectory, and
computational-resource allocation in UAV-aided emergency response
\cite{akter2024satora,akter2025joint}. Although their tasks are computational
rather than relief-delivery tasks, the work supports the broader idea that UAV
task decisions and resource limits should be handled together.

Finally, dynamic disaster settings require some ability to react to new
information. Sayeed et al. and G\"uzey et al. study adaptive UAV coordination
under uncertain or changing disaster conditions
\cite{sayeed2026meanfield,guzey2026fedqmix}. In those works, adaptation is
mainly tied to tracking or localization. PGBM uses a narrower form of
adaptation: a controlled task replacement inside an already planned
relief-delivery mission.

\section{Comparative Analysis and Research Gap}
\label{sec_related_comparison}

Table~\ref{tab:related_summary} compares the most closely related studies
against the modeling dimensions used in PGBM. The purpose of the table is not
to rank the papers, but to show which part of the overall disaster-response
pipeline each work mainly addresses.

\begin{table}[htbp]
\centering
\caption{Comparison of selected UAV disaster-response studies with PGBM.}
\label{tab:related_summary}
\footnotesize
\renewcommand{\arraystretch}{1.15}
\begin{tabularx}{\textwidth}{
    >{\raggedright\arraybackslash}p{2.8cm}
    >{\raggedright\arraybackslash}p{2.5cm}
    >{\raggedright\arraybackslash}p{2.5cm}
    >{\raggedright\arraybackslash}p{3.0cm}
    >{\centering\arraybackslash}X}
\toprule
\textbf{Study} &
\textbf{Main Role} &
\textbf{Routing / 3D Aspect} &
\textbf{Resource Awareness} &
\textbf{Time / Recourse} \\
\midrule
ZaidAlkilani et al. \cite{zaidalkilani2026aiuav}
& Survivor detection and localization
& Partial
& Partial energy consideration
& Not addressed\\
\midrule
Sirma et al. \cite{sirma2025drespnet}
& Disaster-scene understanding
& Partial 3D scene context
& Not a planning model
& Not addressed\\
\midrule
Dorling et al. \cite{dorling2017dronevrp}
& Drone delivery routing
& Delivery route planning
& Payload and battery considered
& Delivery-time limits, no survivor recourse\\
\midrule
Xiong et al. \cite{xiong2023multidrone}
& Mission assignment and 3D path planning
& Explicit 3D path planning
& Partial mission constraints
& Not addressed\\
\midrule
Akter et al. \cite{akter2025joint}
& Trajectory and computation offloading
& Trajectory optimization
& Energy and computation resources
& Latency-aware, not relief recourse\\
\midrule
G\"uzey et al. \cite{guzey2026fedqmix}
& Adaptive victim localization
& Map feasibility considered
& Energy and communication constraints
& Adaptive localization focus\\
\midrule
\textbf{PGBM}
& \textbf{Post-detection relief delivery}
& \textbf{Feasible 3D responder routes}
& \textbf{Payload, energy, reserve, safe return}
& \textbf{Time value and task replacement}\\
\bottomrule
\end{tabularx}
\end{table}

The comparison suggests three connected gaps. First, there is a post-detection
assistance gap. Detection, localization, and scene-understanding studies
provide survivor information, but they usually stop before modeling how
temporary assistance should be delivered while ground rescue is delayed.

Second, there is an integrated resource-aware planning gap. Delivery and
mission-planning studies consider routing, payload, energy, or disaster
logistics, but these elements are usually treated in separate settings. The
reviewed work does not jointly address task selection, 3D responder routing,
changing payload, battery feasibility, and safe return for post-detection
relief delivery.

Third, there is a time-critical and dynamic recourse gap. Existing adaptive UAV
studies mainly focus on changing search, tracking, or localization decisions.
PGBM instead considers whether a newly arriving survivor-assistance task should
replace one uncompleted task in an active relief-delivery mission.

PGBM is therefore positioned as a downstream responder-planning layer. It does
not replace scout-UAV detection, survivor localization, environmental mapping,
or ground-team physical rescue. It uses accepted survivor-assistance tasks as
input and plans how limited responder UAVs should deliver temporary relief
during the critical interval between detection and ground rescue.

\section{Summary}
\label{sec_related_summary}

The reviewed studies provide the sensing, localization, coverage, and
resource-aware planning foundations required for UAV-assisted disaster
response. The remaining need is a compact model for post-detection temporary
assistance, where responder UAVs must choose which tasks to serve, how to route
through a feasible 3D environment, how to respect payload and battery limits,
and how to react when new tasks arrive during an active mission. Chapter~3
develops the proposed PGBM formulation for this downstream planning problem.
